Nanopórové sekvenovanie je obzvlášť užitočné pre jeho schopnosť čítať dlhé sekvencie DNA, čo zjednodušuje zostavovanie genómov a identifikáciu štruktúrnych variácií \cite{Zheng2023Nanopore,Lopez2019DNA}.
Funguje na princípe detekcie zmien v elektrickom prúde, keď molekula prechádza cez nanopór, čo umožňuje priame čítanie sekvencií bez potreby amplifikácie alebo značenia \cite{Zheng2023Nanopore}.
DNA alebo RNA molekuly sú vedené cez nanopór pomocou elektrického poľa, pričom každá báza spôsobuje charakteristickú zmenu prúdu, ktorá je zaznamenaná a analyzovaná softvérom na určenie sekvencie \cite{Zheng2023Nanopore}.
